\subsection{Orders that outlive the trading day}
\label{sec:ha:disabled:long_lived}

Most orders last for the session that placed them, or for the trading day. A
member may also place an order to rest for weeks, and expects to find that order
open each morning without doing anything.

Everything else in this section concerns state that must survive a process or a
machine. These orders must survive \emph{the venue}: every process stopped
deliberately, the day ended, the machines rebooted, and trading started again the
next morning.

\subsubsection{Why the mechanisms already described do not serve}
\label{sec:ha:disabled:long_lived:why}

A checkpoint records the orders held at a point in the sequence. It is useful
only while the record it is anchored to still holds what followed that point, as
\cref{sec:ha:disabled:me:durable} describes. A day boundary breaks that
relationship, because a new day's record does not continue the previous day's.
Retaining every day's record, so that a checkpoint from any day remains usable,
would require the venue to keep every record it has ever written.

An order intended to outlive the day therefore needs a durable store of its own.
The venue holds that store independently of the record of what it did, and reads
it when trading begins rather than recovering it when something fails.

\textbf{No component reaches that store while it is trading.} The venue's
components hold no connection to the store, for the same reason that admission
does not depend on one: where the venue needs the store to be reachable in order
to work, every outage of the store becomes a trading outage. The venue exports
the orders and reads them when trading begins. It brings the store up to date
afterwards from its own record of what it did. The store therefore lags the
venue, and the venue can bring it level again from that record.

That has one consequence for the requirement below. \emph{The venue's own
record} must be durable before the venue tells a member its order has been
cancelled. The store is brought up to date afterwards.

\subsubsection{What the venue must do}
\label{sec:ha:disabled:long_lived:behaviour}

\begin{requirement}{R-0105}{An order that outlives the trading day is held where it survives the venue stopping}
  An order whose terms say it remains open beyond the trading day shall be
  recorded where it survives every component stopping and the day ending.
  \rationale{The venue told the member that the order would rest until the member said
  otherwise. Holding the order only where the day's own state is held makes that
  true for one day, which is not what the member was told.}
  \uncovered{no scenario stops the venue and starts it again with such an order open}
\end{requirement}

\begin{requirement}{R-0106}{Such orders are open when trading next begins}
  When trading begins, every order that was open at the end of the previous day
  and has not since expired shall be open, on the terms it was placed under, and
  cancellable by the member that placed it.
  \rationale{A member that must place its resting orders again each morning does not have
  orders that outlive the day. It has orders that the venue cancels overnight
  without reporting the cancellation.}
  \uncovered{no scenario spans a trading day boundary, which the venue does not mark}
\end{requirement}

\begin{requirement}{R-0114}{The store is known to be readable before trading begins, not during it}
  Whether the store holding these orders can be read shall be established before
  trading starts, and its being unreadable shall stop trading from starting
  rather than be discovered once it has.
  \rationale{The venue reads these orders once, at the start of a day. A store that cannot be
  read at that moment leaves the venue opening with an empty set and nothing to
  indicate that anything is missing. Every member holding such an order finds it
  gone. The check costs little, and the venue must make it before it admits
  members.}
  \uncovered{no scenario starts trading with the store unreadable}
\end{requirement}

\begin{requirement}{R-0107}{A change to such an order is committed before the member is told of it}
  Where such an order is cancelled, amended or filled, that shall be committed to
  the venue's own record before the member is told, and the store the order
  outlives the day in shall be brought up to date from that record.
  \rationale{A member told that its order is cancelled stops watching for that order. Where
  the venue writes nothing durable first, a failure between the two brings the
  order back the next time the venue reads the store. The order is then live, on a
  market that has moved, for a member with no reason to look at it. Requiring the
  venue's own record rather than the store is what allows the store to lag. The
  venue can supply the store from that record with anything the store is missing,
  and it can supply nothing for a change that was never committed anywhere.}
  \uncovered{no scenario cancels such an order and then restarts the venue}
\end{requirement}

\begin{gap}{}
  Orders that outlive the trading day are held only where every other open order
  is held, which does not survive the venue stopping. A member placing one is
  given an order that lasts until the next time the venue is restarted, and is
  told nothing to that effect.
\end{gap}
